% !PRINTED 84
\weberpage{84}
%<<CONT

holds according as \(\lambda\) is different from \(\mu\) or not,
\[
 A x_\mu=0,
\]
and since by our hypothesis \(A\) is different from zero,
\[
 x_\mu=0.
\]

We single out the most frequently used special case \(m=n\), and give the theorem for this case in the following form:

\gesperrt{II. If a system of \(n\) linear homogeneous equations with \(n\) unknowns has a determinant different from zero, then all the unknowns have the value zero; or:}

\gesperrt{If a system of \(n\) linear equations with as many homogeneously occurring unknowns has a solution in which not all the unknowns vanish, then the determinant of the system vanishes.}

Here, by the determinant of a system of \(n\) linear homogeneous equations with \(n\) unknowns is meant the determinant formed from the \(n^2\) coefficients of these equations.

We next consider the case in which \(\nu\) is smaller than \(n\). Since \(m\) must be equal to or greater than \(\nu\), we choose the first \(\nu\) equations of the system (1), and write them thus:
\begin{equation}
\begin{aligned}
 a_1^{(1)}x_1+a_2^{(1)}x_2+\cdots+a_\nu^{(1)}x_\nu&=-a_{\nu+1}^{(1)}x_{\nu+1}-\cdots-a_n^{(1)}x_n,\\
 a_1^{(2)}x_1+a_2^{(2)}x_2+\cdots+a_\nu^{(2)}x_\nu&=-a_{\nu+1}^{(2)}x_{\nu+1}-\cdots-a_n^{(2)}x_n,\\
 &\vdots\\
 a_1^{(\nu)}x_1+a_2^{(\nu)}x_2+\cdots+a_\nu^{(\nu)}x_\nu&=-a_{\nu+1}^{(\nu)}x_{\nu+1}-\cdots-a_n^{(\nu)}x_n.
\end{aligned}
\tag{6}
\end{equation}
We again let \(\mu\) denote one of the indices \(1,2,\ldots,\nu\), multiply equations (6) in succession by \(A_\mu^{(1)},A_\mu^{(2)},\ldots,A_\mu^{(\nu)}\), and add them. It follows, as before, using \S.\,22, (2), (6):
\begin{equation}
 A x_\mu=-C_{\nu+1,\mu}x_{\nu+1}-\cdots-C_{n,\mu}x_n,
\tag{7}
\end{equation}
where, for abbreviation, we have put
\begin{equation}
 C_{\lambda,\mu}=\sum_{1,\nu}^{i}a_\lambda^{(i)}A_\mu^{(i)},
 \qquad \lambda=\nu+1,\nu+2,\ldots,n.
\tag{8}
\end{equation}
By \S.\,22, (2), \(C_{\lambda,\mu}\) is the determinant obtained from the determinant defined by (4) by replacing the
%CONT>>
